\section{The Promise \code{@provides} Makes}
\label{sec:ch11-the-promise-provides-makes}
\index{provides@\code{@provides}|(}%

\code{@requires} says: fetch me something before you call. \code{@provides}
says the opposite. On this path, it says, you do not need to call the owner at
all, because I already have the value.

\index{provides@\code{@provides}!needs a copy of the data}%
That is only ever true when the promising subgraph holds a copy. Mosaic holds
none. Chapter~\ref{ch:the-first-cut-extracting-catalog} gave Catalog its own
database precisely so that no join across the seam would be one connection
string away, and what is left of \code{Product} in \code{Mosaic.Api} is a class
with an identifier on it. There is no title here to promise, no sku, no
category that was not just handed to us by the router.

So Mosaic gets no \code{@provides} in this chapter, and that is the first thing
worth knowing about the directive. A cleanly split system frequently has
nowhere to put one. If you find a place, you have found a copy, and the
question that follows is how the copy stays right.

\subsection*{A subgraph that legitimately has a copy}

\index{samples!crates subgraph}%
The second subgraph in \code{samples/entity-resolution} has one. A
\code{Crate} holds a widget and remembers what that widget was called on the
day it was packed, which is the ordinary reason a service ends up with somebody
else's data: it wrote it down at the time. Two fields return the same object,
and they differ only in what the subgraph promises about it:

\begin{graphqlsdl}[fontsize=\footnotesize]
type Crate {
  widget: Widget! @provides(fields: "name")
  widgetByKey: Widget!
  id: ID!
  label: String!
}

type Widget @key(fields: "id", resolvable: false) {
  id: ID!
  name: String! @external
}
\end{graphqlsdl}

\index{resolvable false@\code{resolvable: false}!the honest use}%
\code{resolvable: false} is the correct declaration here and it is the first
honest use of it in this book. Chapter~\ref{ch:composition} met the flag as the
cause of a satisfiability error and
chapter~\ref{ch:enter-the-router} met it again as a way to break a graph on
purpose. This is what it is for: nothing in \code{crates} can turn a widget key
into a widget, the key exists only so that \code{@provides} has an entity to
point at, and the flag is what stops the router from ever routing one here.

\subsection*{What it buys}

Two plans, from the same router, for two fields that return the same object:

\begin{minted}[fontsize=\footnotesize,breaklines]{text}
widget      @provides(fields: "name")   1 fetch:  crates
widgetByKey no promise                  2 fetches: crates then catalog
the second fetch's path                 crates.@.widgetByKey
\end{minted}

One round trip against two, on a graph where the second round trip is the only
reason the second service is in the plan. That is a real saving and it is the
whole case for the directive.

\index{query plan!path@\code{path}!list traversal}%
The path of that second fetch has something in it that
chapter~\ref{ch:enter-the-router}'s did not. The \code{@} is how the router
writes a step through a list. That chapter saw
\code{browseProducts.nodes}, two object hops in a row, so the notation never
came up.

\subsection*{What it costs}

\index{provides@\code{@provides}!and drift}%
Widget \code{w2} in this sample is packed under a name Catalog does not agree
with. Ask the router for all four crates through both fields and compare what
comes back with what the owner says:

\begin{minted}[fontsize=\footnotesize,breaklines]{text}
Crate one   w1  router says "Hex bolt"        catalog says "Hex bolt"
Crate two   w2  router says "Butterfly nut"   catalog says "Wing nut"
Crate three w3  router says "Split washer"    catalog says "Split washer"
Crate four  w9  router says "Unknown part"    catalog says null
\end{minted}

Three things in one table.

The promise is kept: the router believed \code{crates} and did not check.

The promise is wrong for \code{w2}, and nothing anywhere raises an error. The
client is told a widget is called \enquote{Butterfly nut} by the graph whose
whole purpose is to be the single place that knows what things are called. A
composition check cannot see this, a query plan cannot see this, and no
assertion on the shape of the response can see it. The only thing that can is
knowing that a copy exists.

And the fourth row is the one I did not expect. \code{w9} is a key Catalog has
never heard of, and through the \code{@provides} path it answers perfectly
well. The directive does not only serve stale data, it also conceals a broken
reference: the one query that would have found out is the one the router
decided not to make.

\index{provides@\code{@provides}!when to reach for it}%
Which leaves me somewhere fairly unfashionable on this directive. It is
genuinely the only thing in federation that removes a round trip rather than
arranging one, and I would still reach for it last. The condition for using it
honestly is that your copy is right, and if you can promise that, you have a
cache and you owe it an invalidation story.
Chapter~\ref{ch:resilience-and-performance} is where caching becomes a subject
and that is where the story belongs. What
chapter~\ref{ch:strangling-the-monolith} does with the same tension is
different again: \code{@override} moves a field from one service to another
rather than copying it, which is the answer when the copy was never meant to be
permanent.

\medskip
Every listing in this section came back with data in it. The next one is about
the ones that do not.

\index{provides@\code{@provides}|)}%
